výstup analýzy DNA získaných pomocou nanopórového sekvenovania 
sme si stanovili ako priradenie realneho DNA reťazca ktory bol 
presiel cez nanopore sekvenovanie k jeho zodpovedajúcemu 
Dna nanopore vystupu prestavujuci niekolko sequencií dna, 
ktoré vznikli pri citani dna cez nanopore sekvenovanie.
Ja som sa rozhodol ze namiesto vedenia realneho dna reťazca
hladania jeho respektivneho nanopore vystupu, 
budem sa ho snažit ho vypocitat na zaklade nanopore vystupu,
a to pomocou najdenia dostatočneho mnozstva sekvencií dna v nanopore vystupe,
ktoré su si dostatočne podobne ze mozme predpokladat ze vznikli 
z toho isteho dna reťazca, a na zaklade toho vypocitat konsenzus tychto sekvencií, 
ktory by mal být velmi podobny realnemu dna reťazcu, z ktoreho vznikli.
Dna reťazce su vsak velmi dlhe a da sa predpokladat ze budu obsahovat chyby, 
Co sposobi ze sa budu velmi rozchadzat a budu posunute, preto chceme zistit ich multi alignment,
aby sme robili konsenzus pre tie nukleotidy, ktore sa nachadzaju v rovnakej pozicii.


Vystup by mal teda obsahovat niekolko navzajom zarovnanych sekvencií dna,
a konsenzus tychto zarovnanych sekvencií. A to niekolko krat, pre viacero skupin sekvencií dna, ktore vznikli z toho isteho dna reťazca.
Aby sme mali dostatocne pokritie napriec vsetkym variacami dna katic z dna reťazca.




